\begin{frame}{Coloraciones y progresiones}
    \begin{block}{Definición $k$-coloración}
        Dado un conjunto $X$ una \textit{$r$-coloración} es:
        \begin{itemize}
            \item [(I)] Una función $f:X\to \{1,2,\dots,k\}$
            \item[(II)] Una partición $X=X_{1}\cup X_{2}\cup\dots\cup X_{k}$, donde $X_{i}:=f^{-1}(i)$
        \end{itemize}
        Se dice que $A$ es \textbf{monocromático} si $A\subseteq f^{-1}(i)$ para algún $i$.
    \end{block}

    \begin{block}{Definición progresión aritmética}
        Una \textit{progresión aritmética} es un conjunto de enteros tal que la diferencia de números sucesivos es constante. Es decir, es de la forma
        $$\{a, a+d,a+2d,\dots,a+(k-1)d\}=\{a+id:i=0,1,\dots,k-1\}$$
        Donde $a$ es el inicio, $d$ el paso o distancia y $k$ el tamaño de la progresión.
    \end{block}
\end{frame}

\begin{frame}{Ejemplo de motivación}
    \begin{center}
        ¿Se puede solucionar $x+y=z$ en $\{1,2,3,4,5\}$?
        \pause\\\

        Si, algunas soluciones son:
        \begin{align*}
        1+1 = 2 \qquad
        2+3 = 5 \qquad
        3+1 = 4
        \end{align*}

        \pause\\\

        ¿Se preserva esta propiedad en una 2-coloración de $\{1,2,3,4,5\}$?

        \begin{table}[H]
        \centering
        \begin{tabular}{cccc}
        $\{\}\mid\{1,2,3,4,5\}$ &
        $\{1\}\mid\{2,3,4,5\}$ &
        $\{2\}\mid\{1,3,4,5\}$ &
        $\{3\}\mid\{1,2,4,5\}$ \\

        $\{4\}\mid\{1,2,3,5\}$ &
        $\{5\}\mid\{1,2,3,4\}$ &
        $\{1,2\}\mid\{3,4,5\}$ &
        $\{1,3\}\mid\{2,4,5\}$ \\

        $\{1,4\}\mid\{2,3,5\}$ &
        $\{1,5\}\mid\{2,3,4\}$ &
        $\{2,3\}\mid\{1,4,5\}$ &
        $\{2,4\}\mid\{1,3,5\}$ \\

        $\{2,5\}\mid\{1,3,4\}$ &
        $\{3,4\}\mid\{1,2,5\}$ &
        $\{3,5\}\mid\{1,2,4\}$ &
        $\{4,5\}\mid\{1,2,3\}$ \\
        \end{tabular}
        \end{table}

        \pause

        \textbf{Si se cumple bajo 2-coloración.}
    \end{center}
\end{frame}

\begin{frame}{Teorema de Schur}
    \begin{block}{Teorema/Lema (Schur, 1916)}
        Para todo entero positivo $r$ existe un mínimo entero positivo $S(r)$ tal que toda $r$-coloración de $[n]=\{1,2,\dots,n\}$, con $n\geq S(r)$, contiene una solución monocromática de $x+y=z$.
    \end{block}\pause
    \begin{center}
        \begin{align*}
        S(1)=2 \qquad
        S(2)=5 \qquad
        S(3) = 14 \qquad
        S(4) = 45 \qquad
        S(5) = 161
        \end{align*}    
        \caption{A030126 en OEIS}
    \end{center}

    
    \pause

    \begin{block}{Teorema (Schur, 1916)}
        Para todo natural $k$ existe un entero $n$ tal que para todo primo $p>n$, la ecuación $x^{k}+y^{k}\equiv z^{k}\text{ mod }p$ tiene solución no trivial.
    \end{block}
\end{frame}
